\section{讨论与结论}\label{sec:conclusion}

\emph{主要结论与设计启示：}本文明确了在单发送机约束下将双协议命令共包于单个 \ble{}~5 \texttt{AUX\_ADV\_IND} PDU 何时能在结构、成本与可靠性上获得实际收益。第一，\proto{} 在单包内稳定承载 $26/26$、$52/26$ 与 $28/32$\,B 的 \ble{}/\stdfour{} 完整记录及确定性 $400{+}8b\,\si{\micro\second}$ 包内偏移，在相等单次尝试提交波形时长下将同频每次尝试突发数由 $3$ 减至 $2$（\modejc{} 对 \modesc{}），既不损失计划尝试内联合解析率，也能支撑固定次信道本地时间可用性与单侧断馈恢复。第二，针对无需分片的短命令（$26/26$\,B），传统 \modelc{} 在 $K{=}1$ 时单位成本最低，但在 \SI{6}{ms} 重发预算上限下，\modejc{} 与 \modesc{} 在 $K{=}3$ 即已饱和而 \modelc{} 仍须发射第四次，使单位成本排序发生反转（\modejc{} 低 $14.8\%$、\modesc{} 低 $16.2\%$）；这说明在受限预算策略下，单包大小并非决定累积成本的唯一因素，重发粒度同样关键。第三，当 \ble{} 净荷超过 \SI{27}{B} 直接上限（$52/26$\,B）时，单次信道 PDU 共包免除了五分片 \fragproto{} 的全有或全无重传脆弱性，使 ch3 \modejc{} 在全部三个区组中均高于 ch39 \modelc{}（$K{=}1$ 时 $238/300$ 对 $86/300$；\SI{6}{ms} 时 $291/300$ 对 $93/300$），单位成本降低 $69.4\%\text{--}76.9\%$。

\emph{有效性边界与后续工作：}四项边界制约上述结论的外推：（i）ch3 \modejc{}（\SI{2410}{MHz}）与 ch39 \modelc{}（\SI{2480}{MHz}）对比将格式与载频耦合在一起，$n{=}3$ 区组（$\text{df}{=}2$）下小样本 $t$ 区间跨零，且 B1 点未测量无线单位成本；（ii）成本度量基于 DMA 提交基带时长而非功率计 RF 能耗或切频开销，同时未校准的接收时钟使全部 \num{2730}/\num{3600} 项已联合解析任务的截止期前交付合规不可判定（其余 \num{870}/\num{3600} 项未完成联合解码）；（iii）\stdfour{} 偏移补偿依赖预共享 $b$ 值；（iv）固定次信道 SDR 测试台尚未验证商用 \ble{}/ZigBee 协议栈互操作、动态 \texttt{AuxPtr} 在线时间维持与物理执行器同步。后续工作将引入共享时钟 RF 仪表、同信道对照、商用 \ble{}~5/\stdfour{} 芯片及动态 \texttt{AuxPtr} 跟随以审计严格截止期。
